\chapter{单自由度系统自由振动}
\intro{本章研究单自由度系统在无外力作用下的振动行为，即自由振动。从最简单的无阻尼自由振动出发，深入分析固有频率的物理意义和能量守恒特性;随后讨论各种构型系统的等效刚度和等效质量;最后详细分析有阻尼自由振动的三种振动形态及其判别条件。}

%=====================================================================
\section{无阻尼自由振动}
%=====================================================================

\subsection{运动方程的建立}

考虑一个由质量 $m$ 和弹簧 $k$ 组成的最基本的振动系统，如图 \ref{fig:simple-sdof} 所示。设质量的位移坐标为 $x(t)$，静平衡位置为原点。当质量偏离平衡位置 $x$ 时，弹簧的恢复力为 $-kx$（方向与位移相反）。

根据 Newton 第二定律：
\begin{myformula}
    \[
    m\ddot{x} = -kx
    \]
\end{myformula}

\begin{figure}[H]
\centering
\begin{tikzpicture}[scale=1.2]
    % 墙体
    \fill[vibWall] (0-0.2,0-0.8) rectangle (0+0.2,0+0.8);
    
    % 弹簧
    \draw[decorate,decoration={coil,aspect=0.5,amplitude=6,segment length=8}] (0.5,0) -- (2.5,0) node[midway,above] {$k$};
    
    % 质量块
    \draw[fill=gray!30] (2.5,-0.6) rectangle (3.5,0.6);
    \node at (3.0,0) {$m$};
    
    % 坐标
    \draw[-Stealth] (3.5,0.8) -- (5.0,0.8) node[right] {$x$};
    \draw[thick] (3.5,0.6) -- (3.5,1.0);
    
    % 地面
    \draw[thick] (-0.2,-0.8) -- (5.2,-0.8);
    
    % 支撑
    \draw[fill] (3.0,-0.6) -- (3.0,-0.8) -- (3.05,-0.8) -- (3.05,-0.6);
\end{tikzpicture}
\caption{弹簧-质量系统。}
\label{fig:simple-sdof}
\end{figure}

整理得到标准的运动方程：
\begin{myformula}
    \[
    m\ddot{x} + kx = 0
    \]
\end{myformula}

\subsection{标准形式与固有频率}

将方程两边同时除以质量 $m$，得到：
\begin{myformula}
    \[
    \ddot{x} + \omega_n^2 x = 0
    \]
\end{myformula}

其中：
\begin{myformula}
    \[
    \omega_n = \sqrt{\frac{k}{m}}
    \]
\end{myformula}

$\omega_n$ 称为\textbf{无阻尼固有频率}（Undamped Natural Frequency），单位是 $\text{rad/s}$。它是系统固有的动力学特性参数，仅取决于系统的质量和刚度，与初始条件无关。

固有频率的物理意义：
\begin{itemize}
    \item 系统在无外力作用下自由振动的角频率
    \item 系统刚度和惯性的综合度量：刚度越大（弹簧越"硬"），固有频率越高;质量越大（越"笨重"），固有频率越低
    \item 强迫振动中发生共振的频率条件（无阻尼时 $\omega = \omega_n$）
\end{itemize}

固有频率的常用表达式：
\begin{myformula}
    \[
    f_n = \frac{\omega_n}{2\pi} = \frac{1}{2\pi}\sqrt{\frac{k}{m}}\quad (\text{Hz}),\qquad
    T_n = \frac{1}{f_n} = 2\pi\sqrt{\frac{m}{k}}\quad (\text{s})
    \]
\end{myformula}

\subsection{通解形式}

方程 $\ddot{x} + \omega_n^2 x = 0$ 是二阶常系数线性齐次微分方程，其特征方程为：
\begin{myformula}
    \[
    s^2 + \omega_n^2 = 0 \quad\Rightarrow\quad s = \pm i\omega_n
    \]
\end{myformula}

因此通解为：
\begin{myformula}
    \[
    x(t) = C_1 e^{i\omega_n t} + C_2 e^{-i\omega_n t}
    \]
\end{myformula}

由 Euler 公式，可以写为等效的三角形式：
\begin{myformula}
    \boxed{x(t) = A\cos\omega_n t + B\sin\omega_n t}
    \end{myformula}

或写成幅值-相位形式：
\begin{myformula}
    \boxed{x(t) = X\sin(\omega_n t + \phi)}
    \end{myformula}

两种形式的关系为：
\begin{myformula}
    \[
    X = \sqrt{A^2 + B^2},\qquad \phi = \arctan\frac{B}{A}
    \]
    \[
    A = X\sin\phi,\qquad B = X\cos\phi
    \]
\end{myformula}

\subsection{初始条件求解}

实际振动问题的初始条件通常以初始位移 $x_0$ 和初始速度 $v_0$ 的形式给出：
\[
x(0) = x_0,\qquad \dot{x}(0) = v_0
\]

\textbf{方法一}：使用 $x(t) = A\cos\omega_n t + B\sin\omega_n t$ 形式
\begin{myformula}
    \[
    x(0) = A = x_0,\qquad \dot{x}(0) = \omega_n B = v_0 \;\Rightarrow\; B = \frac{v_0}{\omega_n}
    \]
    \[
    x(t) = x_0\cos\omega_n t + \frac{v_0}{\omega_n}\sin\omega_n t
    \]
\end{myformula}

\textbf{方法二}：使用 $x(t) = X\sin(\omega_n t + \phi)$ 形式
\begin{myformula}
    \[
    X = \sqrt{x_0^2 + \left(\frac{v_0}{\omega_n}\right)^2},\qquad
    \phi = \arctan\frac{\omega_n x_0}{v_0}
    \]
\end{myformula}

\textbf{例 2.1}：已知 $m = 2\,\text{kg}$，$k = 800\,\text{N/m}$，初始条件 $x_0 = 0.02\,\text{m}$，$v_0 = 0\,\text{m/s}$。求系统的自由振动响应。

\textbf{解}：
\[
\omega_n = \sqrt{\frac{k}{m}} = \sqrt{\frac{800}{2}} = 20\,\text{rad/s}
\]
由初始条件：$A = x_0 = 0.02$，$B = v_0/\omega_n = 0$，故
\[
x(t) = 0.02\cos 20t\;\text{m}
\]
周期：$T = 2\pi/\omega_n = 0.314\,\text{s}$，频率：$f_n = \omega_n/(2\pi) = 3.183\,\text{Hz}$。

\subsection{固有频率的物理意义}

固有频率是振动系统最重要的动力学参数之一。从不同角度理解其物理意义：

\begin{itemize}
    \item \textbf{惯性-弹性耦合}：$\omega_n$ 反映了系统在惯性和弹性之间交换能量的速度。$\omega_n$ 越大，能量交换越快，振动越"急促"。
    
    \item \textbf{静态变形法}：设弹簧在重力 $mg$ 下的静变形为 $\delta_{st} = mg/k$，则
    \begin{myformula}
        \[
        \omega_n = \sqrt{\frac{g}{\delta_{st}}}
        \]
    \end{myformula}
    只需测量静变形即可得出固有频率，无需直接知道 $m$ 和 $k$。
    
    \item \textbf{能量观点}：在一个振动周期中，动能和势能以 $2\omega_n$ 的角频率互相转换。在平衡位置，全部能量为动能 $T_{\max} = \frac{1}{2}m\omega_n^2X^2$;在最大位移处，全部能量为势能 $V_{\max} = \frac{1}{2}kX^2$。由 $T_{\max} = V_{\max}$ 得 $\omega_n^2 = k/m$。
\end{itemize}

\subsection{能量守恒}

无阻尼自由振动系统是保守系统，机械能守恒。系统的总能量为：
\begin{myformula}
    \[
    E = T + V = \frac{1}{2}m\dot{x}^2 + \frac{1}{2}kx^2 = \text{const}
    \]
\end{myformula}

代入通解 $x = X\sin(\omega_n t + \phi)$，$\dot{x} = \omega_n X\cos(\omega_n t + \phi)$：
\begin{myformula}
    \[
    E = \frac{1}{2}m\omega_n^2 X^2\cos^2(\omega_n t + \phi) + \frac{1}{2}kX^2\sin^2(\omega_n t + \phi)
      = \frac{1}{2}kX^2\left[\cos^2(\omega_n t + \phi) + \sin^2(\omega_n t + \phi)\right]
      = \frac{1}{2}kX^2
    \]
\end{myformula}

总能量恒为 $\frac{1}{2}kX^2$（或等价地 $\frac{1}{2}m\omega_n^2X^2$），与时间无关。能量在动能和势能之间周期性转换，转换的频率为 $2f_n$。

%=====================================================================
\section{等效刚度与等效质量}
%=====================================================================

实际工程结构往往不是简单的弹簧-质量系统。本节介绍将各种力学构型简化为等效单自由度系统的方法。

\subsection{等效系统法的基本步骤}

\begin{enumerate}
    \item \textbf{确定广义坐标}：选择一个能完整描述系统变形状态的位移坐标
    \item \textbf{计算等效刚度} $k_{eq}$：在广义坐标方向施加单位位移（或广义力），系统弹性恢复力（或力矩）的大小
    \item \textbf{计算等效质量} $m_{eq}$：使等效系统的动能与原系统相等
    \item \textbf{求固有频率}：$\omega_n = \sqrt{k_{eq}/m_{eq}}$
\end{enumerate}

\subsection{悬臂梁端部集中质量}

\begin{center}
\begin{tikzpicture}[scale=1.2]
    % 悬臂梁
    \draw[very thick] (0,0) -- (5,0);
    
    % 固支端
    \draw[thick] (0,0.1) -- (0,-0.3);
    \draw[fill] (-0.1,-0.4) rectangle (0.1,0.2);
    \draw[pattern=north east lines] (-0.1,-0.4) rectangle (0.1,0.2);
    
    % 梁标注
    \node at (2.5,0.3) {$EI, L$};
    
    % 变形后的梁（虚线）
    \draw[thick, dashed] (0,0) .. controls (1.5,-0.1) and (3.5,-0.6) .. (5,-1.5);
    
    % 质量块
    \draw[fill=gray!40] (4.8,-1.5) rectangle (5.4,-1.0);
    \node at (5.1,-1.25) {$m$};
    
    % 位移标注
    \draw[Stealth-Stealth] (5.5,0) -- (5.5,-1.5) node[midway, right] {$x$};
    
    % 力
    \draw[-Stealth, red, thick] (5.1,-1.8) -- (5.1,-2.3) node[right] {$P$};
\end{tikzpicture}
\end{center}

对于端部承受横向力 $P$ 的悬臂梁，根据材料力学：
\begin{myformula}
    \[
    \delta = \frac{PL^3}{3EI}
    \]
\end{myformula}

由刚度定义 $k_{eq} = P/\delta$：
\begin{myformula}
    \boxed{k_{eq} = \frac{3EI}{L^3}}
    \end{myformula}

当梁端附有集中质量 $m$ 时，等效单自由度系统的固有频率为：
\begin{myformula}
    \[
    \omega_n = \sqrt{\frac{k_{eq}}{m}} = \sqrt{\frac{3EI}{mL^3}}
    \]
\end{myformula}

若考虑梁自身分布质量的影响，可用 Rayleigh 法：设梁的振型为 $y(x) = \delta[3(x/L)^2 - (x/L)^3]/2$，则梁的等效质量为 $m_{beam}^* = (33/140)\rho A L \approx 0.2357 m_{beam}$，系统的总等效质量为：
\begin{myformula}
    \[
    m_{eq} = m + \frac{33}{140}\rho A L
    \]
\end{myformula}

\textbf{例 2.2}：钢制悬臂梁长 $L = 0.5\,\text{m}$，截面 $b \times h = 20\times 5\,\text{mm}$，$E = 200\,\text{GPa}$，端部质量 $m = 0.5\,\text{kg}$。求固有频率。

\textbf{解}：截面惯性矩 $I = bh^3/12 = 20\times 5^3/12 = 208.33\,\text{mm}^4 = 2.0833\times 10^{-13}\,\text{m}^4$
\[
k_{eq} = \frac{3EI}{L^3} = \frac{3\times 200\times 10^9\times 2.0833\times 10^{-13}}{0.5^3} = 1000\,\text{N/m}
\]
\[
\omega_n = \sqrt{\frac{1000}{0.5}} = 44.72\,\text{rad/s},\quad f_n = 7.12\,\text{Hz}
\]

\subsection{简支梁中点等效刚度}

对于跨中承受集中力 $P$ 的简支梁，最大挠度为：
\begin{myformula}
    \[
    \delta = \frac{PL^3}{48EI}
    \]
\end{myformula}

因此：
\begin{myformula}
    \boxed{k_{eq} = \frac{48EI}{L^3}}
    \end{myformula}

\subsection{其他常见构型的等效刚度}

\begin{enumerate}
    \item \textbf{两端固支梁中点}：$k_{eq} = \dfrac{192EI}{L^3}$
    \item \textbf{轴向拉伸杆}：$k_{eq} = \dfrac{EA}{L}$
    \item \textbf{扭转圆轴}：$k_t = \dfrac{GJ}{L}$（$J$为极惯性矩）
    \item \textbf{螺旋弹簧}：$k = \dfrac{Gd^4}{64nR^3}$（$d$簧丝直径，$R$弹簧半径，$n$圈数）
\end{enumerate}

\subsection{扭转振动}

对于绕定轴转动的刚体，其扭转振动方程为：
\begin{myformula}
    \[
    J\ddot{\theta} + k_t\theta = 0
    \]
\end{myformula}
其中 $J$ 为绕转轴的转动惯量，$k_t$ 为扭转刚度（$\text{N}\cdot\text{m/rad}$）。

固有频率为：
\begin{myformula}
    \[
    \omega_n = \sqrt{\frac{k_t}{J}}
    \]
\end{myformula}

扭转振动与平动振动的类比关系：
\[
x \leftrightarrow \theta,\quad m \leftrightarrow J,\quad k \leftrightarrow k_t,\quad F \leftrightarrow T
\]

\subsection{能量法求等效参数}

当系统具有分布弹性或分布质量时，使用能量法确定等效参数更为便捷。

\textbf{等效刚度}：由最大势能确定
\begin{myformula}
    \[
    V_{\max} = \frac{1}{2}k_{eq}X^2 \quad\Rightarrow\quad k_{eq} = \frac{2V_{\max}}{X^2}
    \]
\end{myformula}

\textbf{等效质量}：由最大动能确定
\begin{myformula}
    \[
    T_{\max} = \frac{1}{2}m_{eq}\dot{X}_{\max}^2 = \frac{1}{2}m_{eq}(\omega_n X)^2 \quad\Rightarrow\quad m_{eq} = \frac{2T_{\max}}{\omega_n^2 X^2}
    \]
\end{myformula}

\subsection{多个弹簧和质量的等效}

对于由多个弹簧和质量组成的系统，等效刚度和等效质量可以通过系统总势能和总动能来统一计算。

\textbf{例 2.3}：质量 $m$ 由两个弹簧 $k_1$ 和 $k_2$ 对称悬挂，弹簧与铅垂方向夹角为 $\theta$。求系统在铅垂方向微振动的固有频率。

\textbf{解}：设铅垂方向的位移为 $x$（向下为正），则每个弹簧的伸长量为 $x\cos\theta$（小变形假设），弹簧力为 $k_i x\cos\theta$。铅垂方向的总恢复力为 $2kx\cos\theta\cdot\cos\theta = 2kx\cos^2\theta$。故等效刚度为 $k_{eq} = 2k\cos^2\theta$，$\omega_n = \sqrt{2k\cos^2\theta/m}$。

%=====================================================================
\section{有阻尼自由振动}
%=====================================================================

\subsection{运动方程}

引入粘性阻尼力 $F_d = c\dot{x}$ 后，单自由度系统的自由振动方程为：
\begin{myformula}
    \[
    m\ddot{x} + c\dot{x} + kx = 0
    \]
\end{myformula}

\begin{center}
\begin{tikzpicture}[scale=1.2]
    \fill[vibWall] (0-0.2,0-0.8) rectangle (0+0.2,0+0.8);
    
    \draw[decorate,decoration={coil,aspect=0.5,amplitude=5,segment length=7}] (0.5,0) -- (2.0,0) node[midway,above] {$k$};
    
    % 阻尼器
    \draw (2.0,0.4) rectangle (2.8,-0.4);
    \draw (2.0,0.4) -- (2.0,0.8) -- (2.8,0.8) -- (2.8,0.4);
    \draw (2.4,0) -- (2.4,-0.4) node[midway,right] {$c$};
    
    \draw[fill=gray!30] (2.8,-0.5) rectangle (3.6,0.5);
    \node at (3.2,0) {$m$};
    
    \draw[-Stealth] (3.6,0.7) -- (4.8,0.7) node[right] {$x$};
    \draw[thick] (3.6,0.5) -- (3.6,0.9);
    
    \draw[thick] (-0.2,-0.7) -- (5.0,-0.7);
\end{tikzpicture}
\end{center}

\subsection{标准形式与阻尼比}

方程两边除以 $m$：
\begin{myformula}
    \[
    \ddot{x} + \frac{c}{m}\dot{x} + \frac{k}{m}x = 0
    \]
\end{myformula}

引入固有频率 $\omega_n = \sqrt{k/m}$ 和\textbf{阻尼比}（Damping Ratio）$\zeta$：
\begin{myformula}
    \[
    \zeta = \frac{c}{2m\omega_n} = \frac{c}{c_c}
    \]
\end{myformula}

得到标准形式：
\begin{myformula}
    \boxed{\ddot{x} + 2\zeta\omega_n\dot{x} + \omega_n^2 x = 0}
    \end{myformula}

其中 $c_c = 2m\omega_n = 2\sqrt{km}$ 称为\textbf{临界阻尼}（Critical Damping）。$\zeta$ 是无量纲量，用于判断系统的振动形态。

\subsection{特征方程与通解形式}

设解为 $x(t) = Ce^{st}$，代入运动方程得特征方程：
\begin{myformula}
    \[
    s^2 + 2\zeta\omega_n s + \omega_n^2 = 0
    \]
\end{myformula}

特征根为：
\begin{myformula}
    \[
    s_{1,2} = -\zeta\omega_n \pm \omega_n\sqrt{\zeta^2 - 1}
    \]
\end{myformula}

根据 $\zeta$ 的取值，特征根有三种情况，对应三种不同的振动形态。

\subsection{欠阻尼 ($0 < \zeta < 1$)}

当 $\zeta < 1$ 时，$\zeta^2 - 1 < 0$，特征根为共轭复根：
\begin{myformula}
    \[
    s_{1,2} = -\zeta\omega_n \pm i\omega_d
    \]
\end{myformula}
其中 $\omega_d = \omega_n\sqrt{1-\zeta^2}$ 称为\textbf{有阻尼固有频率}（Damped Natural Frequency）。

通解为：
\begin{myformula}
    \[
    x(t) = e^{-\zeta\omega_n t}\left(A\cos\omega_d t + B\sin\omega_d t\right)
    \]
\end{myformula}

或写为幅值-相位形式：
\begin{myformula}
    \boxed{x(t) = Xe^{-\zeta\omega_n t}\sin(\omega_d t + \phi)}
    \end{myformula}

欠阻尼振动的特征：
\begin{itemize}
    \item 振幅按指数规律 $e^{-\zeta\omega_n t}$ 衰减，衰减速率由 $\zeta\omega_n$ 决定
    \item 振动频率 $\omega_d < \omega_n$（阻尼使频率降低）
    \item 当 $\zeta$ 很小时（$\zeta \ll 1$），$\omega_d \approx \omega_n$
    \item 当 $\zeta \to 1$ 时，$\omega_d \to 0$（振动消失）
\end{itemize}

由初始条件 $x(0) = x_0$，$\dot{x}(0) = v_0$ 确定常数：
\begin{myformula}
    \[
    A = x_0,\qquad B = \frac{v_0 + \zeta\omega_n x_0}{\omega_d}
    \]
    \[
    X = \sqrt{x_0^2 + \left(\frac{v_0 + \zeta\omega_n x_0}{\omega_d}\right)^2},\qquad
    \phi = \arctan\frac{\omega_d x_0}{v_0 + \zeta\omega_n x_0}
    \]
\end{myformula}

\subsection{临界阻尼 ($\zeta = 1$)}

当 $\zeta = 1$（即 $c = c_c$）时，特征方程有重根：
\begin{myformula}
    \[
    s_1 = s_2 = -\omega_n
    \]
\end{myformula}

通解为：
\begin{myformula}
    \boxed{x(t) = (A + Bt)e^{-\omega_n t}}
    \end{myformula}

由初始条件：
\begin{myformula}
    \[
    A = x_0,\qquad B = v_0 + \omega_n x_0
    \]
\end{myformula}

临界阻尼是系统刚好不产生振荡的最小阻尼。它将非振荡运动和振荡运动分开，是实际工程设计中（如仪表指针、车辆悬挂系统）经常期望的状态——系统以最快速度返回平衡位置而没有任何超调。

\subsection{过阻尼 ($\zeta > 1$)}

当 $\zeta > 1$（即 $c > c_c$）时，特征根为两个不相等的负实根：
\begin{myformula}
    \[
    s_{1,2} = -\zeta\omega_n \pm \omega_n\sqrt{\zeta^2 - 1}
    \]
\end{myformula}

显然 $s_1 < s_2 < 0$。通解为：
\begin{myformula}
    \boxed{x(t) = Ae^{s_1 t} + Be^{s_2 t}}
    \end{myformula}

由初始条件确定 $A$ 和 $B$：
\begin{myformula}
    \[
    A = \frac{v_0 - s_2 x_0}{s_1 - s_2},\qquad B = \frac{s_1 x_0 - v_0}{s_1 - s_2}
    \]
\end{myformula}

过阻尼系统的运动是非周期的，系统缓慢地返回平衡位置。阻尼越大，$s_2$ 项的衰减越慢，成为主导运动。

\subsection{三种阻尼情况对比}

\begin{center}
\begin{tikzpicture}[scale=1.3]
    % 坐标轴
    \draw[-Stealth] (-0.5,0) -- (7,0) node[right] {$t$};
    \draw[-Stealth] (0,-1.2) -- (0,1.5) node[above] {$x(t)$};
    
    \node[below left] at (0,0) {$O$};
    
    % 欠阻尼：衰减振荡 (zeta=0.1)
    \draw[blue, very thick, domain=0:6.8, samples=200] 
        plot (\x, {exp(-0.15*\x)*cos(deg(1.5*\x))});
    \node[blue] at (1.8,0.6) {欠阻尼 $\zeta=0.1$};
    
    % 临界阻尼 (zeta=1)
    \draw[red, very thick, domain=0:6.8, samples=100] 
        plot (\x, {(1+1.5*\x)*exp(-1.5*\x)});
    \node[red] at (2.5,1.2) {临界阻尼 $\zeta=1$};
    
    % 过阻尼 (zeta=2)
    \draw[green!50!black, very thick, domain=0:6.8, samples=100] 
        plot (\x, {exp(-0.4*\x) + 0.5*exp(-5.6*\x)});
    \node[green!50!black] at (4.5,0.5) {过阻尼 $\zeta=2$};
    
    % 参考线: 包络线
    \draw[blue, dashed] (0,1) -- (6.8,0.36);
    \draw[blue, dashed] (0,-1) -- (6.8,-0.36);
    
    \node at (3.5,-1.0) {$\omega_n = 1.5$};
\end{tikzpicture}
\end{center}

由上图可总结：
\begin{itemize}
    \item \textbf{欠阻尼}：有振荡，振幅逐渐衰减。振荡频率 $\omega_d < \omega_n$
    \item \textbf{临界阻尼}：无振荡，以最快速度返回平衡位置
    \item \textbf{过阻尼}：无振荡，返回平衡位置的速度随 $\zeta$ 增大而变慢
\end{itemize}

\subsection{对数衰减率}

\textbf{对数衰减率}（Logarithmic Decrement）是实验测定阻尼比的重要方法。它定义为相邻两个同方向峰值之比的自然对数。

设 $x_n$ 和 $x_{n+1}$ 分别为 $t_n$ 和 $t_{n+1} = t_n + T_d$ 时刻的两个相邻正峰值（$T_d = 2\pi/\omega_d$ 为有阻尼周期），则：
\begin{myformula}
    \[
    \delta = \ln\frac{x_n}{x_{n+1}}
          = \ln\frac{Xe^{-\zeta\omega_n t_n}}{Xe^{-\zeta\omega_n(t_n + T_d)}}
          = \zeta\omega_n T_d
          = \zeta\omega_n\cdot\frac{2\pi}{\omega_n\sqrt{1-\zeta^2}}
    \]
    \[
    \boxed{\delta = \frac{2\pi\zeta}{\sqrt{1-\zeta^2}}}
    \]
\end{myformula}

当阻尼较小（$\zeta \ll 1$）时，$\sqrt{1-\zeta^2} \approx 1$，近似有：
\begin{myformula}
    \[
    \delta \approx 2\pi\zeta,\qquad \zeta \approx \frac{\delta}{2\pi}
    \]
\end{myformula}

\begin{center}
\begin{tikzpicture}[scale=1.3]
    % 坐标轴
    \draw[-Stealth] (-0.5,0) -- (7.5,0) node[right] {$t$};
    \draw[-Stealth] (0,-1.5) -- (0,1.5) node[above] {$x$};
    \node[below left] at (0,0) {$O$};
    
    % 衰减振荡
    \draw[blue, thick, domain=0:7.2, samples=200] 
        plot (\x, {exp(-0.3*\x)*cos(deg(1.2*pi*\x))});
    
    % 包络线
    \draw[dashed, blue!50] (0,1) -- (7.2,0.115);
    \draw[dashed, blue!50] (0,-1) -- (7.2,-0.115);
    
    % 标注峰值
    \fill[red] (0,1) circle (0.05) node[above right] {$x_1$};
    \fill[red] (1.667,0.5) circle (0.05) node[above] {$x_2$};
    \fill[red] (3.333,0.25) circle (0.05) node[above] {$x_3$};
    \fill[red] (5,0.125) circle (0.05) node[above] {$x_4$};
    
    % 对数衰减标注
    \draw[Stealth-Stealth, red] (0,0.95) -- (1.667,0.95);
    \node[red] at (0.83,1.15) {$T_d$};
    
    \node at (3.5,-0.7) {$\delta = \ln(x_n/x_{n+1})$};
\end{tikzpicture}
\end{center}

若测量相隔 $N$ 个周期的两个峰值 $x_n$ 和 $x_{n+N}$，则：
\begin{myformula}
    \[
    \ln\frac{x_n}{x_{n+N}} = N\delta = \frac{2\pi N\zeta}{\sqrt{1-\zeta^2}}
    \]
\end{myformula}

这在实验测定阻尼比时非常实用——当 $\delta$ 很小时，测量多周期的衰减更能保证精度。

\textbf{例 2.4}（对数衰减率测定阻尼比）：某单自由度系统自由振动实验测得相邻两个振幅分别为 $1.0\,\text{mm}$ 和 $0.8\,\text{mm}$。求系统的阻尼比。

\textbf{解}：
\[
\delta = \ln\frac{1.0}{0.8} = \ln 1.25 = 0.2231
\]
\[
\zeta = \frac{\delta}{\sqrt{4\pi^2 + \delta^2}} = \frac{0.2231}{\sqrt{4\pi^2 + 0.2231^2}} = 0.0355
\]
（近似值 $\zeta \approx \delta/(2\pi) = 0.0355$，误差极小）

\subsection{有阻尼自由振动小结}

\begin{center}
\begin{tabular}{c|c|c|c|c}
\hline
$\zeta$ 范围 & 特征根 & 运动形态 & 衰减方式 & 周期\\
\hline
$\zeta = 0$ & $\pm i\omega_n$ & 等幅简谐振动 & 不衰减 & $T_n = 2\pi/\omega_n$\\
$0<\zeta<1$ & $-\zeta\omega_n\pm i\omega_d$ & 衰减振荡 & 指数衰减 & $T_d = 2\pi/\omega_d$\\
$\zeta = 1$ & $-\omega_n$（重根） & 无振荡，最快返回 & 临界 & $-$\\
$\zeta > 1$ & 两负实根 & 无振荡，缓慢返回 & 两指数衰减 & $-$\\
\hline
\end{tabular}
\end{center}

%=====================================================================
\section{Python 代码示例：有阻尼自由振动}
%=====================================================================

\begin{mycode}{有阻尼自由振动的数值求解}
\begin{lstlisting}[language=Python]
import numpy as np
import matplotlib.pyplot as plt
from scipy.integrate import solve_ivp

# 系统参数
m = 1.0       # kg
k = 100.0     # N/m
omega_n = np.sqrt(k / m)

# 不同阻尼比
zeta_values = [0.05, 0.2, 1.0, 2.0]

# 初始条件
x0 = 0.02   # m
v0 = 0.0    # m/s

# 时间
t_span = (0, 3.0)
t_eval = np.linspace(0, 3, 1000)

fig, ax = plt.subplots(figsize=(10, 6))

for zeta in zeta_values:
    c = 2 * zeta * m * omega_n
    
    def ode_system(t, y):
        x, v = y
        return [v, -(c/m) * v - (k/m) * x]
    
    sol = solve_ivp(ode_system, t_span, [x0, v0],
                    t_eval=t_eval, method='RK45')
    
    label = f'$\\zeta = {zeta}$'
    if zeta == 1.0:
        label += ' (临界阻尼)'
    elif zeta > 1.0:
        label += ' (过阻尼)'
    
    ax.plot(sol.t, sol.y[0], linewidth=1.5, label=label)

ax.set_xlabel('时间 $t$ (s)')
ax.set_ylabel('位移 $x(t)$ (m)')
ax.set_title('不同阻尼比下的自由振动响应')
ax.legend()
ax.grid(True, alpha=0.3)
ax.axhline(y=0, color='k', linewidth=0.5)

plt.tight_layout()
plt.show()

# 对数衰减率计算示例
print("\n=== 对数衰减率计算示例 (zeta=0.05) ===")
zeta = 0.05
omega_d = omega_n * np.sqrt(1 - zeta**2)
delta_exact = 2 * np.pi * zeta / np.sqrt(1 - zeta**2)
delta_approx = 2 * np.pi * zeta
print(f"精确对数衰减率: {delta_exact:.6f}")
print(f"近似对数衰减率: {delta_approx:.6f}")
print(f"相对误差: {abs(delta_exact-delta_approx)/delta_exact*100:.4f}%")
\end{lstlisting}
\end{mycode}
